\begin{lemma}
\label{lemma-product-local}
Any ring with finitely many maximal ideals and
locally nilpotent Jacobson radical is the product of its localizations
at its maximal ideals. Also, all primes are maximal.
\end{lemma}

\begin{proof}
Let the maximal ideals be the distinct ideals
$\mathfrak m_1,\ldots,\mathfrak m_n$, and retain
$I=\bigcap_{i=1}^n\mathfrak m_i$.
If $n=0$, the ring is zero, since every nonzero ring has a
maximal ideal. Its spectrum is empty and its map to the empty
product of rings is the isomorphism of zero rings. Assume $n\geq1$.
Every element of $I$ is nilpotent, so it belongs to every prime.
Since $\mathfrak m_1\cdots\mathfrak m_n\subset I$, a prime
$\mathfrak p$ must contain some $\mathfrak m_i$.
Indeed, otherwise one could choose $a_i\in\mathfrak m_i\setminus
\mathfrak p$ for each $i$, whose product would belong to
$\mathfrak p$ while all its factors avoid it, contradicting
primality. Maximality then gives $\mathfrak p=\mathfrak m_i$.
This proves the assertion about prime ideals.

We construct the precise product isomorphism. For each ordered
pair $i\neq j$, choose $a_{ij}\in\mathfrak m_j$ and
$b_{ij}\in\mathfrak m_i$ with $a_{ij}+b_{ij}=1$, and set
$c_i=\prod_{j\neq i}a_{ij}$. For $n=1$ this product is $1$.
The element $c_i$ has residue $1$ at $\mathfrak m_i$ and
residue $0$ at every other maximal ideal. Thus
$c_i(1-c_i)\in I$. Choose $N_i\geq1$ for which
$[c_i(1-c_i)]^{N_i}=0$. Define the actual element
$$
e_i=\sum_{k=N_i}^{2N_i-1}
\binom{2N_i-1}{k}c_i^k(1-c_i)^{2N_i-1-k}.
$$
All its coefficients and powers are part of this construction.
The binomial identity $(c_i+(1-c_i))^{2N_i-1}=1$ gives
$$
1-e_i=\sum_{k=0}^{N_i-1}
\binom{2N_i-1}{k}c_i^k(1-c_i)^{2N_i-1-k}.
$$
Every term of $e_i$ contains $c_i^{N_i}$, and every term of
$1-e_i$ contains $(1-c_i)^{N_i}$. Their product is zero by
the chosen nilpotence equation. Hence $e_i^2=e_i$.
The residues of $e_i$ at all maximal ideals are the same as
those of $c_i$: substituting $0$ gives zero, and substituting
$1$ leaves just the term with $k=2N_i-1$, equal to one.
For $i\neq j$, the product $e_ie_j$ is an idempotent in $I$.
An idempotent that is nilpotent is zero, since all its positive
powers equal itself. Thus $e_ie_j=0$.
The sum $\sum_i e_i$ is an idempotent, and
$1-\sum_i e_i$ is an idempotent in $I$, so it is zero too.
We have obtained $\sum_i e_i=1$ with pairwise orthogonal
idempotents, without assuming a uniform nilpotence exponent for $I$.

Put $R_i=e_iR$, viewed as a ring with identity $e_i$.
The maps
$$
R\longrightarrow\prod_{i=1}^n R_i,\quad r\longmapsto(e_ir)_i,
\qquad
\prod_{i=1}^n R_i\longrightarrow R,\quad (x_i)_i\longmapsto\sum_i x_i
$$
are inverse ring homomorphisms, by the orthogonality and sum
equations just proved. The first component map is surjective
and has kernel $(1-e_i)R$: $e_ir=0$ implies
$r=(1-e_i)r$, and the reverse implication follows from
$e_i(1-e_i)=0$.
Maximal ideals of $R_i$ correspond to maximal ideals of $R$
containing $(1-e_i)R$. By the computed residues, only
$\mathfrak m_i$ contains $1-e_i$.
Therefore $R_i$ is a nonzero local ring with unique maximal
ideal $e_i\mathfrak m_i$.

The comparison with the original localization is
$$
R_i\longrightarrow R_{\mathfrak m_i},\quad x\longmapsto x/1,
\qquad
R_{\mathfrak m_i}\longrightarrow R_i,\quad
r/s\longmapsto (e_ir)(e_is)^{-1}.
$$
Here $s\notin\mathfrak m_i$ implies
$e_is\notin e_i\mathfrak m_i$: reducing modulo
$\mathfrak m_i$ would otherwise give $s=0$ in its residue
field. Hence $e_is$ is a unit in the original ring $R_i$.
Also $e_i/1$ is an idempotent unit in $R_{\mathfrak m_i}$,
so it equals one there. The first map is consequently unital.
For the second map, a fraction relation
$t(s'r-sr')=0$ with $t\notin\mathfrak m_i$ becomes the
same relation multiplied by $e_i$ in $R_i$; the images of
$t,s,s'$ are units, so the two prescribed images agree.
The maps preserve sums and products by their fraction formulas.
Their composites fix $x\in R_i$ and each original $r/s$,
respectively. Thus both are isomorphisms.
Combining them proves that the canonical map in the assertion is
$$
R\longrightarrow\prod_i R_{\mathfrak m_i},\qquad r\longmapsto(r/1)_i,
$$
with inverse
$(r_i/s_i)_i\mapsto\sum_i(e_ir_i)(e_is_i)^{-1}$.
All terms of this sum are elements of their specified subrings
$e_iR\subset R$; the original denominators and their inverses
in those subrings are retained.

The same construction works for every original $R$-module $M$,
without a finite-generation hypothesis. The maps
$M\to\bigoplus_i e_iM$, $x\mapsto(e_ix)_i$, and
$(x_i)_i\mapsto\sum_i x_i$ are inverse module maps.
The maps $e_iM\to M_{\mathfrak m_i}$, $x\mapsto x/1$,
have inverses
$$
x/s\longmapsto(e_is)^{-1}e_ix.
$$
The same fraction relation and the unit $e_it$ verify that this
inverse is well defined; the composites are identities because
$e_i/1=1$ in the localization.
Consequently $M\to\prod_i M_{\mathfrak m_i}$ is the canonical
isomorphism, with the displayed component maps and inverse sum.
Using finite additivity of length and the surjections
$R\to e_iR$, Lemma \ref{lemma-length-independent} gives
$$
\operatorname{length}_R(M)
=\sum_{i=1}^n\operatorname{length}_{R_{\mathfrak m_i}}
 (M_{\mathfrak m_i}),
$$
also when one or more terms are infinite. In the zero-ring case
every unital module is zero and this is the empty sum.
\end{proof}